\subsection{Adquisición de datos e instancias}

La evaluación del método se apoya en instancias estándar de la biblioteca
\textit{Capacitated Vehicle Routing Problem Library} (CVRPLIB), repositorio ampliamente
utilizado para contrastar algoritmos sobre un conjunto común de casos de prueba y
soluciones de referencia \parencite{cvrplib}. Para permitir una comparación directa con
el método exacto del Informe~1, se reutilizan las mismas cinco instancias de los
Sets~E y~A; adicionalmente, se incorporan cuatro instancias de mayor tamaño que quedaron
fuera del alcance del enfoque exacto, con el fin de evaluar la escalabilidad de PSO.

Cada archivo \texttt{.vrp} se procesa para extraer el número de nodos, la capacidad
vehicular $D$, y la demanda de cada nodo. La mayoría de las instancias sigue la
convención \texttt{EUC\_2D} de TSPLIB utilizada por CVRPLIB, en la que el costo de cada
arista es la distancia euclidiana entre coordenadas, redondeada al entero más cercano,
$c_{ij} = \left[ \sqrt{(x_i-x_j)^2 + (y_i-y_j)^2} \right]$. La instancia \texttt{E-n13-k4}
constituye una excepción: usa el tipo \texttt{EXPLICIT} (formato \texttt{LOWER\_ROW}), en
el que la matriz de distancias se entrega directamente y no existen coordenadas. Dado que
la representación SR-1 (Sección~\ref{sec:representacion}) requiere ubicar puntos de
referencia de vehículos en un plano 2D, para esta instancia se generó un embedding
aproximado mediante \textit{escalamiento multidimensional clásico} (\textit{classical
MDS}) a partir de la matriz de distancias real; esta reconstrucción solo se usa como
insumo interno del mecanismo de referencia del decodificador, mientras que todos los
costos reportados se calculan siempre con la matriz de distancias original de la
instancia, nunca con distancias derivadas del embedding.

La Tabla~\ref{tab:m2_instancias} detalla el conjunto final de instancias evaluadas,
distinguiendo aquellas con óptimo certificado (por el método exacto del Informe~1) de
aquellas para las que solo se dispone de la mejor solución conocida (BKS, \textit{best
known solution}), que no es necesariamente óptima.

\begin{table}[H]
\centering
\caption{Instancias de evaluación. $n$ incluye el depósito; $m$ es el número de
vehículos de la solución de referencia.}
\label{tab:m2_instancias}
\begin{tabular}{lrrrl}
\hline
Instancia & $n$ & $m$ & Referencia & Tipo de referencia \\
\hline
E-n13-k4  & 13  & 4  & 247  & Óptimo certificado (Informe~1) \\
E-n22-k4  & 22  & 4  & 375  & Óptimo certificado (Informe~1) \\
E-n23-k3  & 23  & 3  & 569  & Óptimo certificado (Informe~1) \\
A-n32-k5  & 32  & 5  & 784  & Óptimo certificado (Informe~1) \\
A-n33-k5  & 33  & 5  & 661  & BKS (Informe~1 no certificó: 667 tras 4{,}7~h) \\
\hline
A-n60-k9   & 60  & 9  & 1354 & BKS (CVRPLIB) \\
A-n80-k10  & 80  & 10 & 1763 & BKS (CVRPLIB) \\
E-n76-k10  & 76  & 10 & 830  & BKS (CVRPLIB) \\
E-n101-k8  & 101 & 8  & 815  & BKS (CVRPLIB) \\
\hline
\end{tabular}
\end{table}

Las cuatro instancias bajo la línea separadora quedaron fuera del alcance del método
exacto: extrapolando el crecimiento observado en el Informe~1 (576{,}1~s en $n=32$ a
17\,096{,}4~s, 4{,}7~h, en $n=33$, sin certificar), resolverlas exactamente habría
requerido un tiempo de cómputo no disponible para este trabajo.

\subsection{Diseño del algoritmo PSO}

\subsubsection{Representación: codificación y decodificación}
\label{sec:representacion}

El aspecto central del diseño es la traducción entre el espacio continuo en que opera
PSO y el espacio combinatorio de soluciones del CVRP. Se adoptó la representación
\textbf{SR-1} de \textcite{ai2009particle}, que mantiene la formulación canónica de PSO
sin modificar sus ecuaciones de velocidad/posición (Sección~\ref{sec:representacion} del
Capítulo~3 argumenta esta elección frente a las alternativas de PSO discreto e híbridos).

\paragraph{Codificación.} Para una instancia con $n-1$ clientes y $m$ vehículos, una
partícula es un vector real de $(n-1) + 2m$ dimensiones. Las primeras $n-1$ componentes
son claves de prioridad, una por cliente; las siguientes $2m$ componentes fijan $m$
puntos de referencia $(x_j^{\text{ref}}, y_j^{\text{ref}})$, uno por vehículo, en el
mismo plano donde están los clientes.

\paragraph{Decodificación.} A partir de la partícula se construyen, en orden:
\begin{enumerate}
    \item \textbf{Lista de prioridad de clientes.} Se ordenan las primeras $n-1$
    componentes de menor a mayor valor (\textit{Smallest Position Value}, SPV); el orden
    resultante de los índices de cliente es la lista de prioridad.
    \item \textbf{Matriz de prioridad de vehículos.} Para cada cliente, se ordenan los
    $m$ vehículos de menor a mayor distancia entre el cliente y el punto de referencia de
    cada vehículo.
    \item \textbf{Construcción de rutas.} Se recorre la lista de prioridad de clientes; a
    cada cliente se lo intenta insertar, en la posición de menor costo adicional, en la
    ruta del vehículo de mayor prioridad para ese cliente que aún tenga capacidad
    disponible; si ningún vehículo lo admite sin violar la capacidad, se prueba con el
    siguiente en la prioridad del cliente. Tras cada inserción exitosa se aplica de
    inmediato el procedimiento de mejora local 2-opt sobre esa ruta.
\end{enumerate}

\paragraph{Ejemplo numérico.} Se reproduce el ejemplo original de \textcite{ai2009particle}
para 6 clientes y 2 vehículos. Sea la partícula
\[
\mathbf{x} = (0{,}4,\; 0{,}2,\; 0{,}9,\; 0{,}5,\; 0{,}6,\; 0{,}1 \mid 0{,}8,\; 1{,}5,\; 0{,}3,\; 1{,}2),
\]
donde las primeras 6 componentes son las claves de los clientes $1,\ldots,6$ y las
últimas 4 codifican los puntos de referencia $(0{,}8;\,0{,}3)$ del Vehículo~1 y
$(1{,}5;\,1{,}2)$ del Vehículo~2. Ordenando las claves de los clientes de menor a mayor
($0{,}1 < 0{,}2 < 0{,}4 < 0{,}5 < 0{,}6 < 0{,}9$) se obtiene la lista de prioridad
$6,2,1,4,5,3$. Calculando la distancia de cada cliente a ambos puntos de referencia se
obtiene, por ejemplo, que los clientes $1$ y $2$ están más cerca del Vehículo~2, mientras
que los clientes $3$, $4$, $5$ y $6$ están más cerca del Vehículo~1. Insertando en ese
orden de prioridad, respetando la capacidad, y aplicando 2-opt, se obtienen las rutas
Vehículo~1: $0$-$6$-$3$-$4$-$5$-$0$ y Vehículo~2: $0$-$2$-$1$-$0$.

\paragraph{Factibilidad.} A diferencia de esquemas basados en un \textit{split} sobre una
gran gira (que abren una nueva ruta cada vez que se excede la capacidad), en SR-1 el
número de vehículos $m$ es fijo y viene dado por la partícula misma. Esto implica que,
en casos raros, ningún vehículo puede aceptar a un cliente sin violar su capacidad: el
propio \textcite{ai2009particle} reporta este fenómeno (Tabla~2 de ese trabajo, instancia
\texttt{vrpnc9}: ``\textit{all replications yielded infeasible solution}''). La
Sección~\ref{sec:funcion-aptitud} describe cómo se trata este caso en la implementación.

\subsubsection{Función de aptitud}
\label{sec:funcion-aptitud}

La aptitud de una partícula es el costo total \eqref{eq:m2_objetivo} de las rutas
obtenidas al decodificarla. Cuando la decodificación deja uno o más clientes sin asignar
a ningún vehículo (Sección~\ref{sec:representacion}), se añade una penalización fija y
grande por cada cliente no asignado, de modo que estas partículas queden efectivamente
descartadas por PSO sin impedir su evaluación. No se penaliza el exceso de capacidad
propiamente tal, porque el mecanismo de inserción de SR-1 nunca genera una ruta que
exceda $D$: por construcción, un cliente solo se inserta en un vehículo si su capacidad
lo permite; el único modo de falla es que \emph{ningún} vehículo lo admita, no que se
viole la capacidad de alguno.

\subsubsection{Dinámica del enjambre}

Cada partícula $i$ mantiene una posición $\mathbf{x}_i$ y una velocidad $\mathbf{v}_i$
en $\mathbb{R}^{d}$. En cada iteración, la velocidad y la posición se actualizan según
la formulación canónica con peso de inercia:
\begin{align}
    \mathbf{v}_i(t{+}1) &= w\,\mathbf{v}_i(t)
        + c_1\,\mathbf{r}_1 \otimes \big(\mathbf{p}_i - \mathbf{x}_i(t)\big)
        + c_2\,\mathbf{r}_2 \otimes \big(\mathbf{p}_g - \mathbf{x}_i(t)\big),
        \label{eq:m2_velocidad}\\[2pt]
    \mathbf{x}_i(t{+}1) &= \mathbf{x}_i(t) + \mathbf{v}_i(t{+}1),
        \label{eq:m2_posicion}
\end{align}
donde $\mathbf{p}_i$ es la mejor posición histórica de la partícula (\textit{pbest}),
$\mathbf{p}_g$ la mejor posición global del enjambre (\textit{gbest}, topología global,
sin subdivisión en vecindarios), $w$ el peso de inercia, $c_1$ y $c_2$ los coeficientes
cognitivo y social, $\mathbf{r}_1, \mathbf{r}_2 \sim U(0,1)^d$ vectores aleatorios y
$\otimes$ el producto de Hadamard (componente a componente). El peso de inercia decae
linealmente de $w_0$ a $w_f$ a lo largo de la corrida. La velocidad se acota
componente a componente a $[-v_{\max}, v_{\max}]$; la posición se acota al rango
$[0,1]$ (rango de las claves SPV y de los puntos de referencia), reflejando la velocidad
a $0$ en las componentes que alcanzan el borde. El criterio de término es un número fijo
de iteraciones $T$ (Sección~\ref{sec:parametrizacion}).

\subsubsection{Pseudocódigo}

El Algoritmo~\ref{alg:pso_cvrp} resume el esquema general de PSO para el CVRP, integrando
la codificación, la decodificación a rutas, la evaluación de aptitud y la actualización del
enjambre. La mejora local (2-opt) no se aplica como un paso adicional separado al final de
cada iteración, sino que ya queda incorporada dentro de la propia decodificación
(Sección~\ref{sec:representacion}), inmediatamente después de insertar cada cliente.

\begin{algorithm}[H]
\caption{PSO para el CVRP (representación SR-1)}
\label{alg:pso_cvrp}
\begin{algorithmic}[1]
\Procedure{PSO-CVRP}{instancia, $N$, $w_0$, $w_f$, $c_1$, $c_2$, $v_{\max}$, $T$}
    \State Inicializar posiciones $\mathbf{x}_i \sim U(0,1)^d$ y velocidades $\mathbf{v}_i \gets \mathbf{0}$ del enjambre de $N$ partículas
    \For{cada partícula $i$}
        \State Decodificar $\mathbf{x}_i$ a rutas y evaluar aptitud (Sección~\ref{sec:funcion-aptitud})
        \State $\mathbf{p}_i \gets \mathbf{x}_i$
    \EndFor
    \State $\mathbf{p}_g \gets \arg\min_i \; \text{aptitud}(\mathbf{p}_i)$
    \For{$t \gets 1 \ldots T$}
        \State $w \gets w_0 + \frac{t}{T}(w_f - w_0)$
        \For{cada partícula $i$}
            \State Actualizar $\mathbf{v}_i$ (Ec.~\ref{eq:m2_velocidad}) y $\mathbf{x}_i$ (Ec.~\ref{eq:m2_posicion}); acotar a $v_{\max}$ y a $[0,1]$
            \State Decodificar $\mathbf{x}_i$ y evaluar aptitud
            \State \textbf{Si} aptitud mejora a $\mathbf{p}_i$ \textbf{entonces} $\mathbf{p}_i \gets \mathbf{x}_i$
            \State \textbf{Si} aptitud mejora a $\mathbf{p}_g$ \textbf{entonces} $\mathbf{p}_g \gets \mathbf{x}_i$
        \EndFor
    \EndFor
    \State \textbf{Devolver} rutas decodificadas de $\mathbf{p}_g$
\EndProcedure
\end{algorithmic}
\end{algorithm}

\subsubsection{Parámetros del algoritmo}

Los parámetros del algoritmo y sus valores se resumen en la Tabla~\ref{tab:m2_parametros}.
Su calibración se aborda en la Sección~\ref{sec:parametrizacion}.

\begin{table}[H]
\centering
\caption{Parámetros del algoritmo PSO, encontrados mediante parametrización con Optuna
(Sección~\ref{sec:parametrizacion}).}
\label{tab:m2_parametros}
\begin{tabular}{lll}
\hline
Parámetro & Símbolo & Valor \\
\hline
Tamaño del enjambre       & $N$        & 100 \\
Peso de inercia inicial   & $w_0$      & 0{,}9116 \\
Peso de inercia final     & $w_f$      & 0{,}2782 \\
Coeficiente cognitivo     & $c_1$      & 2{,}4045 \\
Coeficiente social        & $c_2$      & 1{,}1843 \\
Velocidad máxima          & $v_{\max}$ & 0{,}3762 (fracción del rango $[0,1]$) \\
Criterio de parada        & $T$        & 500 iteraciones (derivado del presupuesto, ver más abajo) \\
Repeticiones por instancia& ---        & 31 \\
\hline
\end{tabular}
\end{table}

\subsection{Parametrización}
\label{sec:parametrizacion}

El desempeño de PSO depende de sus parámetros ($N$, $w_0$, $w_f$, $c_1$, $c_2$,
$v_{\max}$), por lo que su calibración se realiza de forma sistemática y reproducible,
separando la fase de \textit{parametrización} de la fase de \textit{evaluación} para
evitar sesgo por sobreajuste: el ajuste se efectúa sobre un subconjunto de instancias de
\textit{tuning} distinto del conjunto de instancias sobre el que luego se reportan los
resultados (Tabla~\ref{tab:m2_instancias}).

\paragraph{Herramienta.} Se usó Optuna \parencite{akiba2019optuna}, con su
muestreador bayesiano por defecto (\textit{Tree-structured Parzen Estimator}, TPE). Se
prefirió Optuna sobre \textit{irace} porque el flujo de trabajo completo (parametrización,
experimentación, análisis) se mantiene íntegramente en Python, sin depender de R como
herramienta externa.

\paragraph{Presupuesto de trials.} Se ejecutaron 60 \textit{trials}, siguiendo la
heurística de fijar un presupuesto de aproximadamente $10\times$ el número de
hiperparámetros a tunear (6 en este caso).

\paragraph{Espacio de búsqueda e instancias de tuning.} La Tabla~\ref{tab:m2_espacio_busqueda}
resume el espacio de búsqueda explorado. Se usaron tres instancias de tuning de tamaños
crecientes y ninguna coincide con las de evaluación: \texttt{A-n34-k5} ($n=34$, chica),
\texttt{A-n62-k8} ($n=62$, mediana) y \texttt{E-n76-k7} ($n=76$, grande). Se usa más de
una instancia, y de tamaños distintos, para que los parámetros encontrados generalicen
razonablemente entre escalas, y no solo a la escala de una única instancia pequeña.

\begin{table}[H]
\centering
\caption{Espacio de búsqueda explorado por Optuna.}
\label{tab:m2_espacio_busqueda}
\begin{tabular}{ll}
\hline
Parámetro & Rango/valores \\
\hline
$N$          & $\{20, 30, 50, 80, 100\}$ \\
$w_0$        & $[0{,}5,\, 0{,}95]$ \\
$w_f$        & $[0{,}1,\, 0{,}5]$ \\
$c_1$        & $[0{,}5,\, 2{,}5]$ \\
$c_2$        & $[0{,}5,\, 2{,}5]$ \\
$v_{\max}$ (fracción del rango) & $[0{,}1,\, 1{,}0]$ \\
\hline
\end{tabular}
\end{table}

\paragraph{Presupuesto de evaluaciones y derivación de $T$.} Cada corrida de PSO dispone
de un presupuesto fijo de $50\,000$ evaluaciones de la función objetivo, igual en la
parametrización y en la experimentación final. Como $N$ es parte del espacio de búsqueda,
el número de iteraciones se deriva de él, $T = \lfloor 50\,000/N \rfloor$, de modo que
todas las configuraciones evaluadas por Optuna consumen el mismo esfuerzo de cómputo
independientemente de cuántas partículas usen.

\paragraph{Resultado.} La mejor configuración encontrada (Tabla~\ref{tab:m2_parametros})
obtuvo un GAP promedio de $8{,}54\%$ sobre las tres instancias de tuning
($3{,}34\%$ en \texttt{A-n34-k5}, $7{,}92\%$ en \texttt{A-n62-k8} y $14{,}37\%$ en
\texttt{E-n76-k7}). El GAP crece con el tamaño de la instancia de tuning, un patrón que
se retoma en la Sección~\ref{sec:m2_comparacion} al discutir cómo generaliza esta
parametrización a las instancias de evaluación más grandes.

\subsection{Implementación computacional}

La implementación se desarrolló en Python 3.12, apoyándose en NumPy
\parencite{harris2020array} para las operaciones vectorizadas del enjambre y en
Numba \parencite{lam2015numba} (compilación \textit{just-in-time}, modo
\texttt{nopython}) para acelerar el núcleo de decodificación SR-1 (construcción de rutas
y 2-opt), que opera con estructuras de largo variable difíciles de vectorizar
directamente con NumPy. Esta aceleración fue necesaria en la práctica: la versión en
Python puro evaluaba entre 680 y 2\,850 partículas por segundo según el tamaño de la
instancia, lo que habría hecho impracticable la parametrización completa con Optuna
(más de 2~horas solo para esa fase); la versión acelerada con Numba, validada para dar
resultados idénticos a la versión de referencia en Python puro, evalúa entre
$130\,000$ y $200\,000$ partículas por segundo, reduciendo la parametrización completa
(60 \textit{trials}) a menos de un minuto. La búsqueda de hiperparámetros se realizó con
Optuna 4.9 \parencite{akiba2019optuna}.

Los experimentos se ejecutaron localmente en un equipo de escritorio con procesador
Intel Core i9-11900K (8 núcleos, 3{,}50\,GHz), 32\,GB de RAM y Windows~11 Home. Este es
el mismo equipo, con la misma configuración de hardware, en el que se ejecutó la
experimentación definitiva del método exacto del Informe~1: si bien las pruebas
iniciales de ese informe se validaron en Google Colab, los tiempos reportados y
utilizados en la comparación de la Sección~\ref{sec:m2_comparacion} (incluidos los
$576{,}1$~s de \texttt{A-n32-k5} y los $17\,096{,}4$~s de \texttt{A-n33-k5}) provienen de
la ejecución definitiva en este mismo equipo local. La comparación de tiempos entre
ambos métodos no está, por lo tanto, sujeta a un sesgo de infraestructura.

% TODO: enlazar el repositorio con el código completo (cvrp_common.py y notebooks de
% parametrización, experimentación y análisis) una vez publicado.

\subsection{Diseño experimental y protocolo de comparación}

La Figura~\ref{fig:m2_diseno_experimental} resume el flujo experimental completo, desde
la parametrización hasta la comparación final con el método exacto.

\begin{figure}[H]
    \centering
    \begin{tikzpicture}[
        box/.style={draw, rounded corners, align=center, minimum height=1.1cm, minimum width=3.1cm, font=\small},
        arr/.style={-{Latex[length=2mm]}, thick},
        node distance=0.9cm and 0.6cm
    ]
        \node[box] (tuning) {3 instancias de tuning\\(A-n34-k5, A-n62-k8, E-n76-k7)};
        \node[box, right=of tuning] (param) {Parametrización\\(Optuna, 60 trials)};
        \node[box, right=of param] (best) {Parámetros óptimos\\(Tabla~\ref{tab:m2_parametros})};
        \node[box, below=of best] (exp) {Experimentación\\(31 repeticiones $\times$ 9 instancias)};
        \node[box, left=of exp] (eval) {9 instancias de evaluación\\(Tabla~\ref{tab:m2_instancias})};
        \node[box, left=of eval] (exacto) {Método exacto\\(Informe~1, 5 instancias)};
        \node[box, below=of exp] (analisis) {Análisis comparativo\\(Sección~\ref{sec:m2_comparacion})};

        \draw[arr] (tuning) -- (param);
        \draw[arr] (param) -- (best);
        \draw[arr] (best) -- (exp);
        \draw[arr] (eval) -- (exp);
        \draw[arr] (exp) -- (analisis);
        \draw[arr] (exacto) -- (analisis);
    \end{tikzpicture}
    \caption{Diseño experimental: las instancias de tuning (parametrización) están
    separadas de las de evaluación (experimentación), y el método exacto del Informe~1
    se incorpora directamente en la etapa de análisis, sin volver a ejecutarse.
    Fuente: elaboración propia.}
    \label{fig:m2_diseno_experimental}
\end{figure}

La \textbf{variable independiente} evaluada es el método de resolución (exacto vs.\ PSO),
contrastado por instancia; dentro de PSO, la semilla de cada repetición es una fuente de
variabilidad estocástica controlada mediante las 31 repeticiones independientes, no una
variable de interés en sí misma. A diferencia del método exacto, que es determinista,
PSO es un algoritmo \textit{estocástico}: distintas semillas producen distintas
trayectorias de búsqueda. Por ello, cada instancia se resuelve con 31 repeticiones
independientes (semillas $0,\ldots,30$), y se reportan estadísticos agregados (mejor,
promedio, desviación estándar) en lugar de un único valor. No se aplica un test
estadístico inferencial en la comparación central de este informe porque uno de los dos
métodos comparados (el exacto) es determinista y entrega un único valor por instancia,
no una distribución: la comparación pertinente es directa, verificando dónde se ubica
ese valor único respecto de la distribución completa de las 31 corridas de PSO, que se
reporta íntegramente (Figura~\ref{fig:m2_boxplot_gap}). Los tests no paramétricos de
comparación de muestras (como Wilcoxon o Friedman) aplican al comparar dos o más
algoritmos estocásticos entre sí, caso que no se da aquí.

Las métricas registradas son: el costo de la mejor solución encontrada y el costo
promedio $\pm$ desviación estándar sobre las 31 repeticiones; el GAP porcentual
respecto al óptimo o BKS de cada instancia,
\begin{equation}
    \text{GAP}\,(\%) = \frac{z_{\text{obtenido}} - z_{\text{referencia}}}{z_{\text{referencia}}} \times 100;
\end{equation}
y el tiempo de ejecución promedio por corrida.

El contraste entre el método exacto y PSO (requisito~8 del enunciado,
Sección~\ref{sec:m2_comparacion}) no iguala el presupuesto de cómputo entre ambos
métodos: mientras que igualar el número de evaluaciones de la función objetivo tendría
sentido al comparar dos metaheurísticas entre sí, no es aplicable aquí, porque el método
exacto no tiene una noción de ``evaluaciones'' comparable y su costo temporal es una
consecuencia directa de certificar (o no) la optimalidad, no un presupuesto que se le
imponga externamente. Por eso, contra el método exacto se reporta el costo propio de
cada enfoque tal como se ejecutó en su respectivo informe.
